\documentclass[11pt]{article}

\usepackage[margin=1in]{geometry}
\usepackage[T1]{fontenc}
\usepackage[utf8]{inputenc}
\usepackage{lmodern}
\usepackage{amsmath,amssymb}
\usepackage{booktabs}
\usepackage{graphicx}
\usepackage{hyperref}
\usepackage{xcolor}
\usepackage{algorithm}
\usepackage{algorithmic}
\usepackage{subcaption}
\usepackage{float}
\usepackage[section]{placeins}


\hypersetup{
  colorlinks=true,
  linkcolor=blue,
  citecolor=blue,
  urlcolor=blue
}

\title{Scalable Self-Play Reinforcement Learning via Hierarchical Spatial Transfer and Curriculum Training: A Case Study on Dots-and-Boxes}
\author{Nishant Mohan}
\date{April 2026}

\begin{document}
\maketitle

\begin{abstract}
Self-play reinforcement learning, exemplified by the AlphaZero paradigm,
has achieved superhuman performance in classical board games such as Go,
Chess, and Shogi. However, applying these methods to combinatorial games
with distinct structural properties---delayed rewards, extra-turn
mechanics, and rapidly growing state spaces---remains underexplored.
We address this gap by developing a self-play system for Dots-and-Boxes,
a deterministic two-player game whose state-space complexity ranges from
$10^5$ ($2{\times}2$) to $10^{25}$ ($5{\times}5$), and whose gameplay
demands tight integration of local tactics with global strategic reasoning.
Our work introduces three methodological contributions:
(i)~a \emph{hierarchical patch-based network} (PatchNet) that transfers
locally learned tactics from small boards to larger ones via frozen
patch-level feature extraction and a trainable global aggregator,
reducing trainable parameters by approximately 60\%;
(ii)~a \emph{phased training curriculum} that co-adapts search budget,
learning rate, and exploration temperature across four stages to
progressively shift the agent from exploration to exploitation;
and (iii)~a \emph{board-size-invariant value target}
($z=\tanh(0.5\,\Delta)$) that preserves gradient magnitude across
board scales, addressing the signal-crushing problem inherent in
score-normalized targets.
Experimental evaluation on $2{\times}2$ and $3{\times}3$ boards shows
that the agent achieves a 100\% win-rate against greedy baselines and
up to 99\% against depth-limited alpha--beta search, while requiring
only hundreds of milliseconds per move compared to minutes for
traditional search at equivalent strategic depth.
\end{abstract}

% ======================================================================
\section{Introduction}
% ======================================================================

Dots-and-Boxes is a deterministic, perfect-information, two-player,
zero-sum pencil-and-paper game. Despite its simple rules---players
alternate drawing edges on a grid; completing the fourth edge of a box
yields a point and an extra turn---the game exhibits deep strategic
complexity. The ``chain rule'' and parity arguments that govern expert
play~\cite{berlekamp2000dots} make it a challenging testbed for
learning-based game agents that must integrate local tactics (capturing
free boxes) with global strategy (managing chain lengths and sacrifices).

The state-space complexity of Dots-and-Boxes grows rapidly: a $3{\times}3$
board already admits $\sim\!10^{10}$ states (Section~\ref{sec:statespace}),
comparable to Connect Four ($\sim\!10^{13}$), while a $5{\times}5$ board
reaches $\sim\!10^{25}$---between Chess ($\sim\!10^{47}$) and Checkers
($\sim\!10^{21}$) in scale~\cite{barker2012solving}. Traditional
alpha--beta search becomes infeasible beyond small boards and modest
depths: even with move ordering and transposition tables, a depth-6
alpha--beta search on a $5{\times}5$ board requires $\sim\!7$~minutes per
move. This motivates learning-based approaches that amortize computation
into a neural network and deploy it within a focused tree search.

We follow the AlphaZero paradigm~\cite{silver2017mastering,silver2018general}:
a policy--value network guides MCTS during self-play; training targets
are derived from search visit counts (policy) and game outcomes (value).
Our contributions are:

\begin{enumerate}
  \item \textbf{PatchNet}: A hierarchical architecture that extracts
    overlapping patches from large boards, evaluates them with a frozen
    pre-trained small-board model, and trains only a lightweight global
    aggregator (Section~\ref{sec:patchnet}).
  \item \textbf{Phased curriculum}: A four-stage training schedule that
    adapts MCTS budget, learning rate, and temperature across training
    (Section~\ref{sec:curriculum}).
  \item \textbf{Board-size-invariant value target}: A scaled-tanh target
    $z=\tanh(0.5\,\Delta)$ that provides consistent gradient magnitude
    across board sizes (Section~\ref{sec:valuetargets}).
  \item \textbf{Systems design}: An efficient training pipeline
    featuring asynchronous MCTS with pending-node scheduling, batched
    neural inference, dihedral symmetry augmentation, and DAG-based
    transposition reuse (Sections~\ref{sec:mcts}--\ref{sec:selfplay}).
\end{enumerate}

% ======================================================================
\section{Background and Related Work}
\label{sec:background}
% ======================================================================

\paragraph{AlphaZero and self-play RL.}
AlphaGo Zero~\cite{silver2017mastering} demonstrated that a single neural
network trained entirely from self-play---without human data---could
surpass human-level performance in Go. The subsequent
AlphaZero~\cite{silver2018general} generalised the approach to Chess and
Shogi, showing that the combination of a policy--value network with MCTS
and the PUCT selection rule~\cite{rosin2011multi} is a powerful
domain-independent framework. MuZero~\cite{schrittwieser2020mastering}
extended this further by learning a latent dynamics model, eliminating the
need for a known game simulator during planning. Our work applies the
AlphaZero paradigm (known rules, no learned dynamics) to
Dots-and-Boxes, a game whose combinatorial structure differs markedly
from Go and Chess.

\paragraph{MCTS foundations.}
Monte Carlo Tree Search was established by Coulom~\cite{coulom2006efficient}
and formalised with regret bounds by the UCT algorithm of Kocsis and
Szepesv\'ari~\cite{kocsis2006bandit}. Modern MCTS implementations use
neural priors (PUCT) and Dirichlet exploration noise at the root.
Our implementation additionally supports asynchronous leaf evaluation
and DAG-based transposition reuse.

\paragraph{Dots-and-Boxes.}
The mathematical theory of Dots-and-Boxes, including chains, loops, and the
Nimstring connection, was developed by
Berlekamp~\cite{berlekamp2000dots}. Barker and
Korf~\cite{barker2012solving} solved $4{\times}5$ Dots-and-Boxes
optimally using enhanced alpha--beta search with transposition tables
and move ordering, establishing computational limits for exact methods.
To our knowledge, no prior work applies deep RL self-play to
Dots-and-Boxes at the scale presented here.

\paragraph{Replay buffers and experience replay.}
Experience replay~\cite{lin1992self} decorrelates training data by
sampling from a buffer of past transitions. We use a capacity-limited
circular buffer with adaptive growth across training phases.

\paragraph{Symmetry in game AI.}
Exploiting board symmetries for data augmentation is standard practice
in AlphaZero-style systems. For Go, the dihedral group $D_4$ provides
8-fold augmentation on square boards~\cite{silver2017mastering}. We
extend this to Dots-and-Boxes, handling the additional complexity of
separate horizontal and vertical edge sets under rotation.

\paragraph{Transfer learning in games.}
Curriculum and transfer approaches have been explored in multi-agent
settings~\cite{bansal2018emergent}. Our PatchNet architecture
(Section~\ref{sec:patchnet}) is a novel form of spatial transfer:
tactical knowledge learned on small boards is reused as a frozen
feature extractor for larger boards.

% ======================================================================
\section{Dots-and-Boxes Formalisation}
\label{sec:game}
% ======================================================================

\subsection{Board Geometry}
A board with $r \times c$ boxes has $(r{+}1)c$ horizontal edges and
$r(c{+}1)$ vertical edges, for an action space of size
\begin{equation}
  |\mathcal{A}| = (r{+}1)c + r(c{+}1) = 2rc + r + c.
\end{equation}
Each action draws one currently undrawn edge. Completing the fourth edge
of a box awards one point and grants an extra turn; otherwise the turn
alternates.

\subsection{Bitboard State Representation}
\label{sec:bitboard}
The game state is encoded using four bitsets and two scalar scores:
\begin{itemize}
  \item \textbf{Horizontal edges} $\mathbf{h} \in \{0,1\}^{(r+1)c}$:
    bit $i$ is set iff horizontal edge $i$ has been drawn.
  \item \textbf{Vertical edges} $\mathbf{v} \in \{0,1\}^{r(c+1)}$:
    analogous for vertical edges.
  \item \textbf{Box ownership} $\mathbf{b}_1, \mathbf{b}_2 \in \{0,1\}^{rc}$:
    bit $j$ is set iff box $j$ is owned by player 1 or 2.
  \item \textbf{Scores} $s_1, s_2 \in \mathbb{Z}_{\ge 0}$: number of
    captured boxes per player.
  \item \textbf{Current player} $p \in \{1, 2\}$.
\end{itemize}
This representation enables $O(1)$ legality checks and incremental
updates via bitwise operations.

\subsection{State-Space Complexity}
\label{sec:statespace}
Let $E = 2rc+r+c$ be the number of edges and $B=rc$ the number of boxes.
An upper bound on the number of distinct game states is
\begin{equation}
  |\mathcal{S}| \le 2^E \cdot 2^B \cdot 2 = 2^{E+B+1},
\end{equation}
where the $2^E$ factor covers edge patterns, $2^B$ covers ownership
assignments for completed boxes\footnote{Each completed box has exactly
two possible owners (the player who drew the fourth edge), so $2^B$ is a
tighter factor than the na\"ive $3^B$ (unclaimed/P1/P2).}, and the final
factor of 2 accounts for the side to move.

\begin{table}[t]
  \centering
  \caption{State-space upper bounds and action-space sizes for representative board sizes.}
  \label{tab:statespace}
  \begin{tabular}{lcccl}
    \toprule
    Board & $E$ & $B$ & $|\mathcal{A}|$ & $|\mathcal{S}|\le$ \\
    \midrule
    $2{\times}2$ & 12 & 4 & 12 & $2^{17} \approx 1.3 \times 10^5$ \\
    $3{\times}3$ & 24 & 9 & 24 & $2^{34} \approx 1.7 \times 10^{10}$ \\
    $4{\times}3$ & 31 & 12 & 31 & $2^{44} \approx 1.8 \times 10^{13}$ \\
    $5{\times}5$ & 60 & 25 & 60 & $2^{86} \approx 7.7 \times 10^{25}$ \\
    $7{\times}8$ & 127 & 56 & 127 & $2^{184} \approx 2.4 \times 10^{55}$ \\
    \bottomrule
  \end{tabular}
\end{table}

Table~\ref{tab:statespace} shows concrete magnitudes. The $5{\times}5$
state space ($\sim\!10^{25}$) lies between Checkers ($\sim\!5 \times
10^{20}$)~\cite{schaeffer2007checkers} and
Chess ($\sim\!10^{47}$)~\cite{silver2018general} in scale, motivating the
use of function approximation and focused tree search rather than
exhaustive computation.

% ======================================================================
\section{Method}
\label{sec:method}
% ======================================================================

\subsection{System Architecture}
\label{sec:architecture}

Figure~\ref{fig:architecture} summarises the training pipeline. Multiple
self-play workers run MCTS in parallel, submitting leaf-evaluation
requests to a centralised batched inference server. Completed games
produce training tuples that are pushed to a replay buffer. A trainer
thread samples minibatches, computes the combined loss, and updates the
network via Adam~\cite{kingma2015adam}. The updated model is
checkpointed and deployed for the next self-play iteration.

\begin{figure}[t]
\centering
\fbox{%
\begin{minipage}{0.95\linewidth}
\small
\begin{center}
\textbf{Training Pipeline Architecture}
\end{center}
\vspace{0.3em}
\begin{tabular}{@{}c@{}}
$W$ \textbf{Self-Play Workers} (parallel threads) \\
$\downarrow$ \\
Each worker: Environment $\xrightarrow{\text{MCTS}}$ action \\
\quad$\hookrightarrow$ leaf nodes $\xrightarrow{\text{submit}}$ \textsc{InferenceServer} \\
\quad$\hookrightarrow$ receive $(\mathbf{p}, v)$, expand, backup \\
\quad$\hookrightarrow$ root visit counts $\Rightarrow$ $\pi$, play action \\
\quad$\hookrightarrow$ on terminal: push $(x, \pi, \text{mask}, z)$ to \textsc{ReplayBuffer}
\end{tabular}
\vspace{0.5em}
\hrule
\vspace{0.5em}
\begin{tabular}{@{}c@{}}
\textsc{InferenceServer} (dedicated thread) \\
accumulate batch ($\le B_{\max}$ states, 200\,$\mu$s window) \\
$\rightarrow$ single forward pass $\rightarrow$ softmax $\rightarrow$ dispatch $(\mathbf{p}, v)$
\end{tabular}
\vspace{0.5em}
\hrule
\vspace{0.5em}
\begin{tabular}{@{}c@{}}
\textbf{Trainer} (main thread) \\
sample minibatches from \textsc{ReplayBuffer} \\
$\rightarrow$ compute $\mathcal{L} = \mathcal{L}_\pi + 1.5\,\mathcal{L}_v$ \\
$\rightarrow$ gradient clipping ($\|\nabla\|_2 \le 1.0$) \\
$\rightarrow$ Adam step $\rightarrow$ checkpoint $\theta$
\end{tabular}
\end{minipage}%
}
\caption{End-to-end architecture of the training pipeline.
Self-play workers, the inference server, and the trainer operate concurrently.}
\label{fig:architecture}
\end{figure}

% ------------------------------------------------------------------
\subsection{Policy--Value Network (AlphaZeroBitNet)}
\label{sec:network}
% ------------------------------------------------------------------

The base model is a fully-connected residual network with a shared trunk
and two output heads.

\paragraph{Input features.}
The state is encoded as a flat feature vector composed of:
\begin{enumerate}
  \item \textbf{Horizontal edge bits} $\in \{0,1\}^{(r+1)c}$
  \item \textbf{Vertical edge bits} $\in \{0,1\}^{r(c+1)}$
  \item \textbf{Box ownership} $\in \{-1,0,+1\}^{rc}$: $+1$ if owned by
    the current player, $-1$ if by the opponent, $0$ if unclaimed
  \item \textbf{Edge-count one-hot} $\in \{0,1\}^{rc \times 5}$: for each
    box, a 5-channel one-hot encoding of the number of drawn surrounding
    edges (0--4). The count-3 feature directly signals ``free box'' (one
    move from capture), while count-2 signals a half-open chain
  \item \textbf{Game-phase scalar}: $\text{boxes\_filled} / B \in [0,1]$
  \item \textbf{Score-advantage scalar}: $(s_{\text{me}} - s_{\text{opp}}) / B$
\end{enumerate}
Yielding a total input size of
$(r{+}1)c + r(c{+}1) + rc + 5rc + 2 = 2rc + r + c + 6rc + 2$.

\paragraph{Architecture.}
\begin{itemize}
  \item \textbf{Input projection:}
    $\text{Linear}(d_{\text{in}}, H) \rightarrow \text{LayerNorm}
    \rightarrow \text{ReLU}$.
  \item \textbf{Residual tower:} $N_B$ fully-connected residual blocks.
    Each block applies $\text{Linear} \rightarrow \text{LayerNorm}
    \rightarrow \text{ReLU} \rightarrow \text{Dropout}(p) \rightarrow
    \text{Linear} \rightarrow \text{LayerNorm}$, then adds a skip
    connection and applies ReLU. LayerNorm~\cite{ba2016layer} is used
    instead of BatchNorm since self-play minibatches can be small.
  \item \textbf{Policy head:}
    $\text{Linear}(H, H/2) \rightarrow \text{LayerNorm} \rightarrow
    \text{ReLU} \rightarrow \text{Linear}(H/2, |\mathcal{A}|)$,
    producing logits over all edges.
  \item \textbf{Value head (deep MLP):}
    $\text{Linear}(H, H/2) \rightarrow \text{LayerNorm} \rightarrow
    \text{ReLU} \rightarrow \text{Linear}(H/2, H/4) \rightarrow
    \text{LayerNorm} \rightarrow \text{ReLU} \rightarrow
    \text{Linear}(H/4, 1) \rightarrow \tanh$.
\end{itemize}
The value head is deliberately deeper than the single-layer head used in
standard AlphaZero~\cite{silver2018general}: Dots-and-Boxes requires
aggregating global signals such as total score advantage, game phase, and
chain parity, which benefit from additional nonlinear capacity.

\paragraph{Illegal move masking.}
During both training and inference, an additive mask is applied before
the softmax:
\begin{equation}
  \hat{\mathbf{p}} = \text{softmax}\!\Big(\mathbf{l} + (\mathbf{m} - \mathbf{1}) \cdot 10^9\Big),
\end{equation}
where $\mathbf{l}$ are the raw logits and
$\mathbf{m} \in \{0,1\}^{|\mathcal{A}|}$ is the legal-action mask. The
large negative bias drives illegal-action probabilities to zero without
requiring architecture changes. The replay buffer stores this mask
alongside each training sample to ensure consistent masking during
gradient computation.

% ------------------------------------------------------------------
\subsection{Hierarchical Patch-Based Network (PatchNet)}
\label{sec:patchnet}
% ------------------------------------------------------------------

Training a monolithic network on larger boards (e.g., $5{\times}5$,
$7{\times}8$) is expensive because both the network capacity and MCTS
budget must scale with the state-space size. We introduce PatchNet, a
hierarchical architecture that transfers locally learned tactics from a
small pre-trained model to larger boards.

\paragraph{Core idea.}
Local tactical patterns in Dots-and-Boxes---such as ``avoid giving away
a free box'' (count-3 avoidance) or ``capture a chain''---are
\emph{spatially local}: they depend only on a small neighbourhood of the
board. PatchNet exploits this by:
\begin{enumerate}
  \item Extracting all overlapping $k{\times}k$ patches from the
    $r{\times}c$ board (yielding
    $(r{-}k{+}1)(c{-}k{+}1)$ patches).
  \item Running a \emph{frozen} pre-trained $k{\times}k$ AlphaZeroBitNet
    on each patch, producing local policy logits and a local value signal.
  \item Scattering local policy logits to global action indices
    (averaging where patches overlap).
  \item Concatenating the global feature representation, the scattered
    policy map, and the per-patch values into a fusion vector.
  \item Processing the fusion vector through a trainable global
    aggregator (residual tower + policy/value heads) to produce the final
    output.
\end{enumerate}

\paragraph{Architecture details.}
For a $5{\times}5$ board with $3{\times}3$ patches, there are
$3 \times 3 = 9$ overlapping patches. Each patch produces 24 action
logits and 1 value signal. The global aggregator has its own residual
tower ($N_G$ blocks, hidden dimension $H_G$). The fusion input dimension
is $H_G + |\mathcal{A}| + P$, where $P$ is the number of patches.

\paragraph{Transfer training protocol.}
\begin{enumerate}
  \item \textbf{Pre-train} a standard AlphaZeroBitNet on the $k{\times}k$
    board until convergence.
  \item \textbf{Freeze} the local model: all parameters have
    \texttt{requires\_grad=false} and the model runs in evaluation mode
    (no dropout).
  \item \textbf{Train only the global aggregator} on the large board via
    normal self-play. Since local tactics are pre-learned, the bootstrap
    phase can be skipped entirely.
\end{enumerate}

\paragraph{Advantages.}
\begin{itemize}
  \item $\sim$60\% fewer trainable parameters than a monolithic network
    of comparable capacity.
  \item Faster convergence: local tactical patterns do not need
    re-discovery.
  \item Natural spatial decomposition without weight-sharing assumptions
    of CNNs.
\end{itemize}

This approach is analogous to CNN local receptive fields but operates on
the graph structure of the Dots-and-Boxes grid rather than a regular
pixel lattice.

% ------------------------------------------------------------------
\subsection{Monte Carlo Tree Search}
\label{sec:mcts}
% ------------------------------------------------------------------

We implement AlphaZero-style MCTS with the PUCT selection
rule~\cite{rosin2011multi}:
\begin{equation}
  a^* = \arg\max_a \left(
    Q(s,a) + c_{\text{puct}} \, P(s,a) \,
    \frac{\sqrt{\textstyle\sum_b N(s,b)}}{1 + N(s,a)}
  \right),
\end{equation}
where $P(s,a)$ is the policy prior from the network, $N(s,a)$ are visit
counts, and $Q(s,a) = W(s,a)/N(s,a)$ is the mean action value.

\paragraph{Dirichlet exploration noise.}
At the root node during self-play, priors are perturbed:
\begin{equation}
  P'(s,a) = (1-\varepsilon)\,P(s,a) + \varepsilon\,\eta_a,
  \quad \eta \sim \text{Dir}(\alpha).
\end{equation}

\paragraph{First-play urgency (FPU) reduction.}
Unvisited actions ($N(s,a)=0$) use a pessimistic value estimate:
\begin{equation}
  Q_{\text{FPU}}(s,a) = \bar{Q}(s) - \delta_{\text{fpu}},
\end{equation}
where $\bar{Q}(s)$ is the parent's mean value and
$\delta_{\text{fpu}} \in [0, 1]$ is the FPU reduction parameter. This
prevents premature uniform exploration when the policy prior is
uncalibrated early in training.

\paragraph{Extra-turn value conversion.}
Dots-and-Boxes includes ``extra move'' captures where the current player
retains the turn. During PUCT selection, if a child node has the
\emph{same} player as the parent (an extra-turn edge), $Q(s,a)$ is used
directly; if the player changes, $Q(s,a)$ is \emph{negated} to convert
to the parent's perspective. This is critical for correctly valuing
capture sequences.

\paragraph{Temperature-based action selection.}
The action played in self-play is sampled proportionally to visit counts
raised to a temperature:
\begin{equation}
  \pi(a) \propto N(s,a)^{1/\tau}, \qquad
  \tau = \begin{cases}
    \tau_{\text{explore}} & \text{if move count} < t_{\text{thresh}}, \\
    \tau_{\text{exploit}} & \text{otherwise}.
  \end{cases}
\end{equation}
This schedule encourages diverse openings (high $\tau$) while playing
near-greedily ($\tau \to 0$) in the mid- and endgame.

% ------------------------------------------------------------------
\subsection{Asynchronous MCTS}
\label{sec:async}
% ------------------------------------------------------------------

In synchronous MCTS, each simulation blocks at every leaf until the
neural network returns $(\mathbf{p}, v)$. When inference is the
bottleneck, this underutilises CPU cores. We implement an asynchronous
variant using \emph{pending nodes}:

\begin{algorithm}[t]
\caption{Asynchronous MCTS root search (conceptual)}
\label{alg:async}
\begin{algorithmic}[1]
\FOR{$\text{sim} = 1$ \TO $S$}
  \STATE $n \gets \text{root}$
  \WHILE{$n$ is expanded \AND not terminal}
    \STATE $n \gets \text{SelectChild}(n)$ via PUCT
    \IF{$n.\text{status} = \text{PENDING}$}
      \STATE \textbf{continue} to next simulation
    \ENDIF
  \ENDWHILE
  \IF{$n$ is terminal}
    \STATE \text{Backpropagate}(path, outcome)
  \ELSE
    \STATE $n.\text{status} \gets \text{PENDING}$
    \STATE \text{EnqueueInference}($n.\text{state}$)
  \ENDIF
\ENDFOR
\STATE \textbf{On inference reply} ($n$, $\mathbf{p}$, $v$):
\STATE \quad $\text{Expand}(n, \mathbf{p})$; \text{Backpropagate}(path, $v$)
\end{algorithmic}
\end{algorithm}

When a simulation reaches an unexpanded leaf, it submits an inference
request, marks the node \textsc{Pending}, and immediately starts a new
simulation. If a simulation encounters a pending node, it skips that path.
When the inference result arrives, the node transitions to
\textsc{Expanded} with initialised priors. This overlaps CPU-side
selection/backup with inference and increases simulations per second.

% ------------------------------------------------------------------
\subsection{Tree vs.\ DAG Search}
\label{sec:dag}
% ------------------------------------------------------------------

The implementation supports both tree-based search and a hybrid DAG mode
controlled by a configuration flag:

\begin{itemize}
  \item \textbf{Training (self-play):} Tree mode (default). Each
    simulation path allocates fresh nodes, avoiding shared-statistics
    effects during learning while policies and values are unstable.
  \item \textbf{Evaluation:} DAG mode. A per-move transposition table
    maps a 128-bit hash of the full board state to shared node pointers.
    Identical states reached via different move orders reuse the same
    node, reducing duplicated expansions. The hash uses golden-ratio
    mixing of four bitset hash pairs plus the current player, providing
    collision probability $< 2^{-64}$ for boards up to $20{\times}20$.
    The table is cleared after each move to bound memory.
\end{itemize}

Additionally, an \emph{inference cache} avoids redundant neural network
evaluations within a single \texttt{act()} call: if the same board state
is reached by multiple simulation paths, the cached policy--value result
is reused without re-querying the inference server.

% ------------------------------------------------------------------
\subsection{Capture Boost Prior}
\label{sec:captureboost}
% ------------------------------------------------------------------

During the early bootstrap phase when the network's policy is essentially
random, MCTS may fail to discover that completing a box is always
beneficial. To address this cold-start problem, we optionally apply a
\emph{capture boost}: the network's prior probability for actions that
complete at least one box is multiplicatively boosted by a factor
$\beta_{\text{cap}}$ before normalisation. We linearly decay
$\beta_{\text{cap}}$ to zero across training phases as the network
learns to prioritise captures on its own.

% ------------------------------------------------------------------
\subsection{Training Pipeline}
\label{sec:selfplay}
% ------------------------------------------------------------------

\subsubsection{Training Targets and Loss}
\label{sec:loss}

For each state $s_t$ visited during self-play, we store a tuple
$(x_t, \pi_t, \mathbf{m}_t, z)$ where $x_t$ is the encoded state,
$\pi_t$ is the policy target from MCTS visit counts, $\mathbf{m}_t$ is
the legal-action mask, and $z$ is the value target.

The network output is $(\hat{\mathbf{p}}, \hat{v}) = f_\theta(x)$. We
minimise the loss:
\begin{equation}
  \mathcal{L}(\theta) =
    \underbrace{-\sum_a \pi(a) \log \hat{p}(a)}_{\mathcal{L}_\pi}
    + 1.5 \underbrace{(\hat{v} - z)^2}_{\mathcal{L}_v}
  \label{eq:loss}
\end{equation}
where the policy loss uses masked log-softmax
($\hat{p}(a) = \text{softmax}(\mathbf{l} + (\mathbf{m}-1)\cdot 10^9)_a$).
The value loss coefficient of 1.5 (rather than the standard 1.0 used in
AlphaZero~\cite{silver2018general}) was found empirically to stabilise
early training by providing a stronger gradient signal for value
estimation.

\paragraph{Optimisation.}
We use Adam~\cite{kingma2015adam} with per-phase learning rates.
Gradient norms are clipped to $\|\nabla\|_2 \le 1.0$ to prevent
instability from rare outlier states. The number of gradient steps per
epoch is capped at $\min(|\text{buffer}|/B_{\text{batch}},\, 256)$ to
ensure stable updates even as the replay buffer grows.

% ------------------------------------------------------------------
\subsubsection{Value Target Design}
\label{sec:valuetargets}
% ------------------------------------------------------------------

Dots-and-Boxes yields final scores $s_1, s_2$ with margin
$\Delta = s_{\text{me}} - s_{\text{opp}}$. We implement five value
targets, offering a spectrum from sparse to dense supervision:

\begin{table}[t]
\centering
\caption{Value target definitions. $B = rc$ is the total number of boxes,
$\Delta = s_{\text{me}} - s_{\text{opp}}$.}
\label{tab:valuetargets}
\begin{tabular}{llc}
\toprule
Name & Formula & Range \\
\midrule
\textsc{WinLoss} & $z = \text{sign}(\Delta)$ & $\{-1,0,+1\}$ \\
\textsc{ScoreDiff} & $z = \Delta / B$ & $[-1, +1]$ \\
\textsc{ScoreDiffSqrt} & $z = \text{sign}(\Delta)\sqrt{|\Delta|/B}$ & $[-1, +1]$ \\
\textsc{ScoreDiffTanh} & $z = \tanh(2\Delta / B)$ & $(-1, +1)$ \\
\textsc{ScoreDiffScaled} & $z = \tanh(0.5 \cdot \Delta)$ & $(-1, +1)$ \\
\bottomrule
\end{tabular}
\end{table}

\paragraph{The signal-crushing problem.}
The \textsc{ScoreDiffTanh} target divides by $B$ before applying $\tanh$.
On large boards ($B = 25$ for $5{\times}5$), a 1-box lead yields
$z = \tanh(2/25) = \tanh(0.08) \approx 0.08$, providing minimal gradient
signal. Values cluster in $[-0.2, +0.2]$, making it difficult for the
network to distinguish winning from losing positions.

\paragraph{Board-size-invariant target.}
\textsc{ScoreDiffScaled} addresses this by applying $\tanh$ to the
\emph{raw} (unnormalized) score difference:
$z = \tanh(0.5 \cdot \Delta)$. This gives:
\begin{center}
\begin{tabular}{ccccc}
$\Delta=1$ & $\Delta=2$ & $\Delta=3$ & $\Delta=5$ \\
$z \approx 0.46$ & $z \approx 0.76$ & $z \approx 0.91$ & $z \approx 0.99$
\end{tabular}
\end{center}
The mapping is independent of board size: a 1-box lead always yields
$z \approx 0.46$ whether on a $3{\times}3$ or $7{\times}8$ board,
providing a strong, discriminative training signal.

% ------------------------------------------------------------------
\subsubsection{Phased Training Curriculum}
\label{sec:curriculum}
% ------------------------------------------------------------------

Rather than using fixed hyperparameters throughout training, we employ a
multi-phase curriculum that adapts key parameters as the agent improves.
Each phase specifies its own MCTS simulation budget, learning rate,
temperature schedule, and capture boost:

\begin{table}[t]
\centering
\caption{Phased training curriculum (example for $5{\times}5$ board).
$S$: MCTS sims, $\eta$: learning rate, $\tau_e/\tau_x$: explore/exploit
temperature, $\beta$: capture boost.}
\label{tab:phases}
\begin{tabular}{lcccccl}
\toprule
Phase & Iters & $S$ & $\eta$ & $\tau_e/\tau_x$ & $\beta$ & Purpose \\
\midrule
Bootstrap   & 30 & 250  & 0.003    & 1.0/0.3  & 0.3 & Cold-start; learn captures \\
ChainAware  & 80 & 600  & 0.001    & 0.9/0.15 & 0.0 & Chain reasoning \\
DeepSearch  & 70 & 1000 & 0.0003   & 0.8/0.08 & 0.0 & Endgame/parity \\
Mastery     & 50 & 1400 & 0.0001   & 0.6/0.02 & 0.0 & Near-greedy exploitation \\
\bottomrule
\end{tabular}
\end{table}

The \textbf{Bootstrap} phase uses lower simulation budgets and higher exploration
to rapidly accumulate diverse experience. The optional capture boost ensures
MCTS always explores capture moves during the cold-start period.
\textbf{ChainAware} increases MCTS depth so the agent can discover
multi-step chain strategies. \textbf{DeepSearch} further increases
simulations to refine endgame and parity play. \textbf{Mastery} uses
near-greedy exploitation with very low learning rate for final polishing.

For PatchNet, the Bootstrap phase can be skipped entirely since local
tactics are pre-learned from the frozen small-board model.

% ------------------------------------------------------------------
\subsubsection{Self-Play, Replay Buffer, and Batched Inference}
\label{sec:replay}
% ------------------------------------------------------------------

\paragraph{Self-play workers.}
Training is parallelised across $W$ worker threads. Each worker runs
complete self-play games using MCTS and pushes training tuples to a
shared replay buffer. Workers communicate with the inference server via
asynchronous submit/poll: the worker submits a state snapshot and polls
for the response without blocking.

\paragraph{Batched inference server.}
A dedicated inference thread aggregates requests from all workers into
batches of up to $B_{\max}$ states. After a brief accumulation window
($\sim$200\,$\mu$s) or when the batch reaches capacity, the server
runs a single forward pass, applies softmax to obtain policy probabilities,
and dispatches results back to workers. This amortises GPU/CPU kernel
launch overhead and improves hardware utilisation.

\paragraph{Replay buffer.}
The replay buffer is a capacity-limited circular queue implemented as a
thread-safe \texttt{std::deque}. When the buffer is full, the oldest
samples are discarded (FIFO eviction). The buffer capacity grows
adaptively across iterations (e.g., +2000 per iteration) to accommodate
the expanding diversity of self-play experience. Random minibatches are
drawn using a partial Fisher--Yates shuffle for
$O(\text{batch\_size})$ sampling complexity.

For reproducibility and training resumption, the buffer supports
serialisation: state, policy, mask tensors and value scalars are stacked
and written to a LibTorch archive. Optimizer state and phase/iteration
progress are checkpointed alongside the model weights.

% ------------------------------------------------------------------
\subsubsection{Symmetry Augmentation}
\label{sec:symmetry}
% ------------------------------------------------------------------

Dots-and-Boxes boards admit geometric symmetries that preserve the game
dynamics. We exploit these for zero-cost data augmentation by remapping
state features and policy targets under symmetry transformations.

\begin{itemize}
  \item \textbf{Rectangular boards} ($r \ne c$): horizontal flip and
    vertical flip (2 transforms $\Rightarrow$ 3$\times$ data with identity).
  \item \textbf{Square boards} ($r = c$): the full dihedral group $D_4$
    comprising identity, 2 reflections, 3 rotations, and 2 diagonal
    flips (7 transforms $\Rightarrow$ 8$\times$ data).
\end{itemize}

Unlike image augmentation where pixel indices transform uniformly,
Dots-and-Boxes requires separate handling of horizontal and vertical
edge sets: under rotation, horizontal edges map to vertical edges and
vice versa. The augmentation module pre-computes index permutation
tables for each transform, then generates augmented training samples
by bitset remapping---a pure $O(|E|)$ operation with no neural network
calls.

Currently, symmetry augmentation is disabled for PatchNet because
patch features would require separate permutation tables; extending
augmentation to PatchNet is left for future work.

% ======================================================================
\section{Implementation Details}
\label{sec:implementation}
% ======================================================================

\paragraph{Language and framework.}
The system is written in C++17 and uses the PyTorch C++ frontend
(LibTorch) for neural network definition, training, and inference.
CUDA is used automatically when available; the system degrades
gracefully to CPU-only execution.

\paragraph{Hyperparameters.}
Table~\ref{tab:hyperparams} lists the default hyperparameters for each
board size.

\begin{table}[t]
\centering
\caption{Default network and MCTS hyperparameters by board size.}
\label{tab:hyperparams}
\begin{tabular}{lcccc}
\toprule
Parameter & $2{\times}2$ & $3{\times}3$ & $5{\times}5$ & $5{\times}5$ (Patch) \\
\midrule
Hidden size $H$ & 64 & 128 & 384 & 256 (global) \\
Residual blocks $N_B$ & 4 & 6 & 10 & 6 (global) \\
Dropout $p$ & 0.1 & 0.1 & 0.06 & 0.08 \\
MCTS sims $S$ & 200 & 400--800 & 600--1400 & 400--1200 \\
$c_{\text{puct}}$ & 1.5 & 1.6 & 1.6 & 1.6 \\
Dirichlet $\alpha$ & 0.3 & 0.15 & 0.10 & 0.10 \\
Dirichlet $\varepsilon$ & 0.25 & 0.20 & 0.25 & 0.25 \\
FPU reduction & 0.25 & 0.20 & 0.20 & 0.20 \\
Workers $W$ & 8 & 16 & 16 & 16 \\
Batch size & 128 & 256 & 256 & 256 \\
Buffer capacity & 20K & 50K & 60K & 50K \\
Value target & WL & SD Tanh & SD Scaled & SD Scaled \\
Augmentation & $D_4$ & $D_4$ & off & off \\
\bottomrule
\end{tabular}
\end{table}

\paragraph{Reproducibility.}
The full source code, including all configuration files and evaluation
scripts, is publicly available at
\url{https://github.com/Nishant040305/Dot-Boxes}.

% ======================================================================
\section{Experimental Setup}
\label{sec:experiments}
% ======================================================================

\subsection{Baselines}
We evaluate against the following opponents:
\begin{itemize}
  \item \textbf{Random}: uniform random legal move selection.
  \item \textbf{Greedy}: always captures a free box (count-3) if
    available; otherwise plays randomly. This is the minimal tactical
    baseline.
  \item \textbf{Alpha--Beta}: depth-limited minimax with alpha--beta
    pruning, transposition tables, and move ordering. Depths tested:
    $d \in \{2, 3, 4, 6\}$ where computationally feasible.
\end{itemize}

\subsection{Evaluation Protocol}
For each agent-vs-baseline pair, we play $n \ge 100$ evaluation games
with alternating first player. The agent uses greedy action selection
(temperature $\tau \to 0$) during evaluation. We report:
\begin{itemize}
  \item \textbf{Win-rate} (\% of games won by the trained agent).
  \item \textbf{Average score difference} ($\bar{\Delta}$).
  \item \textbf{Compute budget}: MCTS simulations per move and wall-clock
    time per move.
\end{itemize}

\subsection{Board Sizes}
\begin{itemize}
  \item \textbf{$2{\times}2$}: 12 edges. Rapid iteration and debugging.
  \item \textbf{$3{\times}3$}: 24 edges. Standard benchmark where chains and
    tactical patterns emerge.
  \item \textbf{$5{\times}5$}: 60 edges. Tests scalability; PatchNet vs.\ 
    monolithic comparison.
\end{itemize}

% ======================================================================
\section{Results}
\label{sec:results}
% ======================================================================

\subsection{Win-Rates and the Effect of Simulation Budget}

We evaluate the agent at varying MCTS simulations per move to understand how planning scales with computation. Performance against stronger opponents (Greedy, Alpha--Beta) generally improves with increased simulations, indicating that the agent effectively leverages deeper search to outmaneuver tactical play. On the $3{\times}3$ board, the agent achieves a 100.0\% win-rate against Greedy with 400 simulations per move and maintains this performance at higher budgets. Against a depth-3 Alpha--Beta opponent, the agent wins 95.0\% of games with 800 simulations, improving further to 99.0\% at 3200 simulations.

\paragraph{Non-monotonic scaling on $2{\times}2$.}
A notable exception to the above trend is observed on the $2{\times}2$ board, where win-rates against weaker opponents (Random, Greedy) exhibit a \emph{non-monotonic, approximately parabolic relationship} with the simulation budget. Specifically, increasing simulations from 200 to 800 leads to a decrease in win-rate, followed by a recovery at higher budgets (e.g., 3200 simulations).

\paragraph{GTO vs.\ exploitative behavior.}
This phenomenon arises from the distinction between \emph{game-theoretic optimal (GTO)} play and \emph{exploitative} play. With larger simulation budgets, MCTS increasingly evaluates positions under the assumption of a perfect opponent and therefore prioritizes \emph{risk-averse, worst-case optimal strategies}. On the $2{\times}2$ board—where many lines provably lead to draws—the agent prefers moves that guarantee a draw over moves that could yield a win only if the opponent errs. As a result, the agent effectively \emph{trades off exploitability for optimality}, leading to increased draw rates and reduced empirical win-rate against suboptimal opponents.

\paragraph{Effect of low simulation budgets.}
At lower simulation counts, the agent relies more heavily on its learned policy prior, resulting in more diverse and less constrained play. This induces more complex and unstable board states, which increases the likelihood of mistakes from weaker opponents. Consequently, the agent achieves higher win-rates despite being theoretically less optimal.

\paragraph{Second-mover advantage in $2{\times}2$.}
Additionally, the $2{\times}2$ environment exhibits a strong \emph{second-mover advantage}. Due to the small state space, optimal play from the second player often yields a forced win or decisive advantage, whereas the first player is largely restricted to defensive play that frequently results in draws. This structural asymmetry further amplifies sensitivity to simulation depth and contributes to the observed non-monotonic (parabolic) scaling behavior.
\begin{table}[H]
\centering
\caption{Comprehensive percentage of Wins (W), Losses (L), and Draws (D) of the trained AlphaZero agent against baselines at varying MCTS simulation budgets.}
\label{tab:winrates_sims}
\begin{tabular}{llcccccc}
\toprule
Board & Opponent & Sims/move & W (\%) & L (\%) & D (\%) & $\bar{\Delta}$ & Time/move \\
\midrule
$2{\times}2$ & Random & 200  & 87.5 & 0.0 & 12.5 & +2.90 & 88.6 ms \\
$2{\times}2$ & Random & 400  & 85.2 & 0.0 & 14.8 & +2.79 & 140.7 ms \\
$2{\times}2$ & Random & 800  & 83.8 & 0.0 & 16.2 & +2.88 & 162.5 ms \\
$2{\times}2$ & Random & 1600 & 84.2 & 0.2 & 15.5 & +2.86 & 275.0 ms \\
$2{\times}2$ & Random & 3200 & 91.8 & 0.0 & 8.2 & +2.99 & 323.0 ms \\
\midrule
$2{\times}2$ & Greedy & 200  & 54.8 & 2.8 & 42.5 & +1.79 & 87.4 ms \\
$2{\times}2$ & Greedy & 400  & 52.8 & 2.2 & 45.0 & +1.83 & 143.3 ms \\
$2{\times}2$ & Greedy & 800  & 58.5 & 1.2 & 40.2 & +1.88 & 205.9 ms \\
$2{\times}2$ & Greedy & 1600 & 63.2 & 1.0 & 35.8 & +2.00 & 241.7 ms \\
$2{\times}2$ & Greedy & 3200 & 76.0 & 3.8 & 20.2 & +2.28 & 331.2 ms \\
\midrule
$2{\times}2$ & Alpha--Beta ($d{=}3$) & 200  & 50.0 & 15.0 & 35.0 & +1.21 & 73.5 ms \\
$2{\times}2$ & Alpha--Beta ($d{=}3$) & 400  & 50.0 & 12.8 & 37.2 & +1.29 & 110.2 ms \\
$2{\times}2$ & Alpha--Beta ($d{=}3$) & 800  & 51.2 & 12.2 & 36.5 & +1.39 & 104.3 ms \\
$2{\times}2$ & Alpha--Beta ($d{=}3$) & 1600 & 50.0 & 12.0 & 38.0 & +1.33 & 233.5 ms \\
$2{\times}2$ & Alpha--Beta ($d{=}3$) & 3200 & 70.8 & 5.2 & 24.0 & +2.31 & 257.4 ms \\
\midrule
$3{\times}3$ & Random & 200  & 100.0 & 0.0 & 0.0 & +8.08 & 282.9 ms \\
$3{\times}3$ & Random & 400  & 100.0 & 0.0 & 0.0 & +8.11 & 472.1 ms \\
$3{\times}3$ & Random & 800  & 100.0 & 0.0 & 0.0 & +8.12 & 679.4 ms \\
$3{\times}3$ & Random & 1600 & 100.0 & 0.0 & 0.0 & +8.24 & 1040.3 ms \\
$3{\times}3$ & Random & 3200 & 100.0 & 0.0 & 0.0 & +8.33 & 2110.8 ms \\
\midrule
$3{\times}3$ & Greedy & 200  & 100.0 & 0.0 & 0.0 & +5.41 & 273.4 ms \\
$3{\times}3$ & Greedy & 400  & 100.0 & 0.0 & 0.0 & +5.46 & 430.0 ms \\
$3{\times}3$ & Greedy & 800  & 100.0 & 0.0 & 0.0 & +5.43 & 748.0 ms \\
$3{\times}3$ & Greedy & 1600 & 100.0 & 0.0 & 0.0 & +5.52 & 1545.0 ms \\
$3{\times}3$ & Greedy & 3200 & 100.0 & 0.0 & 0.0 & +5.83 & 2880.0 ms \\
\midrule
$3{\times}3$ & Alpha--Beta ($d{=}3$) & 200  & 94.0 & 6.0 & 0.0 & +3.40 & 358.8 ms \\
$3{\times}3$ & Alpha--Beta ($d{=}3$) & 400  & 96.0 & 4.0 & 0.0 & +3.34 & 367.2 ms \\
$3{\times}3$ & Alpha--Beta ($d{=}3$) & 800  & 95.0 & 5.0 & 0.0 & +3.42 & 584.1 ms \\
$3{\times}3$ & Alpha--Beta ($d{=}3$) & 1600 & 100.0 & 0.0 & 0.0 & +3.90 & 946.5 ms \\
$3{\times}3$ & Alpha--Beta ($d{=}3$) & 3200 & 99.0 & 1.0 & 0.0 & +3.68 & 1810.2 ms \\
\bottomrule
\end{tabular}
\vspace{0.3em}

\noindent\textit{Note: $5{\times}5$ board experiments are ongoing and will be reported in an extended version of this work.}
\end{table}

\subsection{Depth-Based Analysis Against Alpha--Beta}

We further challenge the agent against the Alpha--Beta search with varying analytical depths. For this comparison, the AlphaZero agent plays with a fixed 3200 simulations per move.

\begin{table}[H]
\centering
\caption{AlphaZero (3200 sims/move) vs Alpha--Beta search at increasing depths shown as normalized percentages.}
\label{tab:depth_analysis}
\begin{tabular}{lcccccc}
\toprule
Board & AB Depth & W (\%) & L (\%) & D (\%) & $\bar{\Delta}$ & Time/move \\
\midrule
$2{\times}2$ & $d{=}4$ & 62.0 & 12.0 & 26.0 & +1.80 & 330.4 ms \\
$2{\times}2$ & $d{=}5$ & 58.0 & 22.0 & 20.0 & +0.72 & 287.2 ms \\
$2{\times}2$ & $d{=}6$ & 50.0 & 28.0 & 22.0 & +0.44 & 282.2 ms \\
$2{\times}2$ & $d{=}7$ & 50.0 & 22.0 & 28.0 & +0.56 & 293.1 ms \\
$2{\times}2$ & $d{=}8$ & 50.0 & 36.0 & 14.0 & +0.28 & 312.7 ms \\
$2{\times}2$ & $d{=}9$ & 50.0 & 50.0 & 0.0 & +0.00 & 324.5 ms \\
$2{\times}2$ & $d{=}10$ & 50.0 & 50.0 & 0.0 & +0.00 & 307.3 ms \\
\midrule
$3{\times}3$ & $d{=}3$ & 99.0 & 1.0 & 0.0 & +3.68 & 1810.2 ms \\
$3{\times}3$ & $d{=}4$ & 96.0 & 4.0 & 0.0 & +3.42 & 1852.6 ms \\
$3{\times}3$ & $d{=}5$ & 94.0 & 6.0 & 0.0 & +3.28 & 1984.7 ms \\
$3{\times}3$ & $d{=}6$ & 92.0 & 8.0 & 0.0 & +2.84 & 2568.7 ms \\
$3{\times}3$ & $d{=}7$ & 92.0 & 8.0 & 0.0 & +2.76 & 2570.2 ms \\
\bottomrule
\end{tabular}
\end{table}

Despite playing against deep alpha--beta search capable of seeing up to $d=10$ plies ahead on the $2{\times}2$ board and $d=7$ plies ahead on the $3{\times}3$ board, the AlphaZero agent maintains a strong advantage. On the $3{\times}3$ board, it achieves over a 90\% win rate even against $d{=}7$. On the $2{\times}2$ board, it maintains at least a 50\% win rate against exhaustive search (up to $d{=}10$).

\subsection{Scalability: AlphaZero MCTS vs.\ Alpha--Beta}

\begin{table}[t]
\centering
\caption{Compute comparison between alpha--beta and AlphaZero MCTS.
Alpha--beta times are from sequential search on an Intel i7-13700H.}
\label{tab:scalability}
\begin{tabular}{llcc}
\toprule
Method & Board & Depth/Sims & Time/move \\
\midrule
Minimax         & $5{\times}5$ & $d{=}3$ & $\sim$2 min \\
Alpha--Beta     & $5{\times}5$ & $d{=}3$ & $\sim$30 s \\
Alpha--Beta     & $5{\times}5$ & $d{=}4$ & $\sim$1 min \\
Alpha--Beta     & $5{\times}5$ & $d{=}6$ & $\sim$7 min \\
\bottomrule
\end{tabular}
\end{table}

Traditional search methods exhibit exponential time growth with depth.
Alpha--beta pruning and move ordering improve feasibility (depth-6
reduced from infeasible minimax to $\sim$7~minutes), but remain
impractical for real-time play on boards larger than $5{\times}5$.
In contrast, the AlphaZero agent amortises planning computation into the
neural network, maintaining bounded per-move latency (typically under
2~seconds on the $3{\times}3$ board even at 3200 simulations) regardless
of the underlying state-space complexity.

\subsection{PatchNet vs.\ Monolithic}

The PatchNet architecture is designed to reduce the computational burden
of training on larger boards by reusing locally learned features. On the
$5{\times}5$ board, PatchNet employs a frozen $3{\times}3$ pre-trained
model to extract features from 9 overlapping patches, training only a
lightweight global aggregator (256 hidden units, 6 residual blocks)
compared to the monolithic network (384 hidden units, 10 residual blocks).
This yields an estimated $\sim$60\% reduction in trainable parameters.
Full comparative evaluation on the $5{\times}5$ board---including
convergence speed, final win-rates, and parameter counts---is ongoing
and will be reported in a future extended version.

\subsection{Training Dynamics}

Across all board sizes, we observe a consistent training progression.
During the Bootstrap phase, the policy loss decreases rapidly as the
agent learns basic capture tactics. The transition to the ChainAware
phase---marked by increased MCTS simulations and reduced
exploration---correlates with the emergence of multi-step chain
reasoning, visible as a temporary increase in value loss followed by
convergence. On the $3{\times}3$ board, the agent achieves 100\%
win-rate against the Greedy baseline within the first two training
phases, confirming that local tactical mastery is acquired early.
Subsequent phases refine positional judgement and endgame play,
leading to strong performance against deeper alpha--beta search.

% ======================================================================
\section{Ablation Studies}
\label{sec:ablations}
% ======================================================================

Ablation experiments are designed to isolate the contribution of individual
components. The following table outlines the planned comparisons; full
results will be included upon completion of the experimental campaign.

\begin{table}[t]
\centering
\caption{Ablation results on the $3{\times}3$ board. Each row modifies one
component from the full configuration. Win-rate is measured against Greedy
over 200 games.}
\label{tab:ablations}
\begin{tabular}{lcc}
\toprule
Configuration & Win-rate (\%) & $\Delta$ from full \\
\midrule
Full system (baseline)              & -- & -- \\
No Dirichlet noise ($\varepsilon=0$)                 & -- & -- \\
No edge-count one-hot features                        & -- & -- \\
WinLoss value target (instead of ScoreDiffTanh)       & -- & -- \\
No symmetry augmentation                              & -- & -- \\
FPU reduction $= 0$                                    & -- & -- \\
Half MCTS simulations (200 instead of 400)            & -- & -- \\
Double MCTS simulations (800)                         & -- & -- \\
\bottomrule
\end{tabular}
\end{table}

% ======================================================================
\section{Analysis and Discussion}
\label{sec:discussion}
% ======================================================================

\paragraph{Dominance over heuristic baselines.}
On the $3{\times}3$ board, the agent achieves a perfect 100\% win-rate
against both Random and Greedy baselines across all tested simulation
budgets (200--3200), with average score differences of $+8.3$ and $+5.8$
respectively. This indicates that the agent has fully internalised both
capture tactics and deeper strategic reasoning (e.g., chain management),
rendering heuristic opponents non-competitive.

\paragraph{Scaling of search budget.}
Increasing the MCTS simulation budget generally improves performance
against stronger opponents. On the $3{\times}3$ board, win-rate against
depth-3 alpha--beta search improves from 94\% at 200 simulations to
100\% at 1600 simulations. However, on the $2{\times}2$ board, we
observe a non-monotonic relationship: higher simulation counts initially
reduce win-rate against weak opponents before recovering at very high
budgets. We attribute this to the agent converging toward
game-theoretically optimal (GTO) play---preferring guaranteed draws over
risky wins---when given sufficient planning depth in the small, highly
constrained state space.

\paragraph{Robustness against deep search.}
Even against alpha--beta search with increasing depth, the agent remains
competitive. On the $3{\times}3$ board, it maintains a 92\% win-rate
against $d{=}7$ alpha--beta with 3200 simulations per move. On the
$2{\times}2$ board, the agent achieves 50\% win-rate against $d{=}9$
and $d{=}10$ (exhaustive search), reflecting convergence to optimal
play where wins are evenly split by starting position. The graceful
degradation pattern (from 70.8\% at $d{=}3$ to 50\% at $d{=}10$)
suggests the agent approaches near-optimal play.

\paragraph{Computational efficiency.}
The AlphaZero agent computes moves in hundreds of milliseconds to low
single-digit seconds, whereas alpha--beta search on the $5{\times}5$
board requires 30~seconds to 7~minutes per move even at moderate depths.
This orders-of-magnitude improvement in per-move latency demonstrates
the practical advantage of amortising strategic knowledge into a neural
network.

\paragraph{GTO--exploitative trade-off.}
The non-monotonic scaling pattern on the $2{\times}2$ board highlights
a fundamental tension in self-play systems: agents trained with deeper
search converge toward minimax-optimal strategies that may underperform
against suboptimal opponents. This finding has implications for
deployment, where opponents are rarely optimal, and suggests that
adaptive simulation budgets or opponent-modelling may be beneficial.

% ======================================================================
\section{Limitations and Future Work}
\label{sec:limitations}
% ======================================================================

\paragraph{Current limitations.}
\begin{itemize}
  \item Dots-and-Boxes has known strategic motifs (chains, parity,
    nimstring values~\cite{berlekamp2000dots}) that the fully-connected
    network must learn from scratch. Explicit chain-detection features or
    graph neural networks could provide useful inductive bias.
  \item PatchNet currently does not support symmetry augmentation because
    patch features require separate permutation tables.
  \item The fully-connected architecture does not exploit the spatial
    locality of the board as efficiently as convolutional or
    graph-based architectures might.
  \item Training was conducted only on CPU; GPU acceleration would
    enable scaling to larger boards and higher simulation budgets.
\end{itemize}

\paragraph{Future directions.}
\begin{itemize}
  \item \textbf{Convolutional and graph-based architectures}: replacing
    the FC backbone with a GNN or CNN that respects the grid topology.
  \item \textbf{Curriculum learning across board sizes}: progressively
    increasing board size during training, using PatchNet-style transfer.
  \item \textbf{Learned dynamics model} (MuZero-style~\cite{schrittwieser2020mastering}):
    eliminating the need for an explicit game simulator during planning.
  \item \textbf{Evaluation against human experts} and optimal solvers
    on small boards~\cite{barker2012solving}.
  \item \textbf{Multi-agent population training} to avoid strategy
    collapse from pure self-play.
\end{itemize}

% ======================================================================
\section{Conclusion}
\label{sec:conclusion}
% ======================================================================

We presented a self-play reinforcement learning system for
Dots-and-Boxes that extends the AlphaZero paradigm with three
methodological contributions: the PatchNet hierarchical architecture for
spatial transfer of locally learned tactics, a four-phase training
curriculum that co-adapts search depth and exploration across training,
and a board-size-invariant value target that avoids gradient signal
crushing on larger boards.

On the $3{\times}3$ board, the trained agent achieves a 100\% win-rate
against greedy baselines and up to 99\% against depth-3 alpha--beta
search, while maintaining 92\% even against depth-7 search. The agent
computes moves in under 2~seconds compared to minutes required by
traditional alpha--beta at comparable strategic depth. On the
$2{\times}2$ board, analysis of the non-monotonic scaling behaviour
reveals a trade-off between game-theoretically optimal and exploitative
play, providing insight into the interaction between search depth and
opponent strength.

These results demonstrate that AlphaZero-style learning is effective for
combinatorial games with delayed rewards and extra-turn mechanics, and
that hierarchical spatial transfer offers a promising direction for
scaling learned game-playing agents to larger and more complex domains.

\bibliographystyle{plain}
\bibliography{refs}

\end{document}
